\begin{trapbox}
	\begin{description}[style=nextline, leftmargin=2.8em, labelsep=0.5em, itemsep=0.35em]
		\item[\textbf{\textcolor{CTrap}{易错点一（忽视定义域）}}]直接使用切化弦时忘记要求$\cos x\neq0$，当$\cos x=0$时仍机械同除，导致错误。
		\item[\textbf{\textcolor{CTrap}{正确理解（检查前提）：}}]用公式前必须确认$\cos x\neq0$；若$\cos x=0$，需单独讨论原式的值。
		\item[\textbf{\textcolor{CTrap}{错因分析（条件遗漏）：}}]学生只记公式形式，忽略公式成立的必要条件“分母不为零且$\cos x\neq0$”。
	\end{description}

	\medskip
	\begin{description}[style=nextline, leftmargin=2.8em, labelsep=0.5em, itemsep=0.35em]
		\item[\textbf{\textcolor{CTrap}{易错点二（齐次误判）}}]分式并非齐次（各项次数不同）时，错误地使用同除$\cos^n x$，导致变形不合法。
		\item[\textbf{\textcolor{CTrap}{正确理解（验证次数）：}}]使用前必须验证分子分母每一项中$\sin x$与$\cos x$的总次数是否相同。
		\item[\textbf{\textcolor{CTrap}{错因分析（概念混淆）：}}]将部分项次数相同误认为齐次，缺乏整体检查。
	\end{description}

	\medskip
	\begin{description}[style=nextline, leftmargin=2.8em, labelsep=0.5em, itemsep=0.35em]
		\item[\textbf{\textcolor{CTrap}{易错点三（分母零忽略）}}]只担心$\cos x=0$，忽略原分母$Q_n(\sin x,\cos x)$可能为零的点，导致定义域扩大。
		\item[\textbf{\textcolor{CTrap}{正确理解（全面考虑）：}}]使用切化弦后，原式的定义域不变，应保证原分母不为零且$\cos x\neq0$。
		\item[\textbf{\textcolor{CTrap}{错因分析（范围扩大）：}}]化简后的有理式可能简化了限制，学生误代入原来被排除的$x$值。
	\end{description}
\end{trapbox}
